% This is text associated with the open textbook "Open Electromagnetics"
% See immediately below for author, date
% This work is licensed under the Creative Commons Attribution-ShareAlike 4.0 International License. (CC BY-SA 4.0) 
% To view a copy of the license, visit http://creativecommons.org/licenses/by-sa/4.0/

%ORB TITLE Reflection Coefficient for Various Terminations
%ORB AUTHOR 0        # S.W. Ellingson, Virginia Tech (ellingson@vt.edu)
%ORB DATE 20180802   # yyyymmdd
%ORB ABSTRACT Voltage reflection coefficients for various terminations 

\label{m0085_Reflection_Coefficient_for_Various_Terminations} % <-- Must match name of this file and module directory name.  Don't mess with this!

We now consider values of $\Gamma$  
that arise for commonly-encountered terminations.

{\bf Matched Load ($Z_L=Z_0$).} 
\index{load!matched}
In this case, the termination may be a device with impedance $Z_0$, or the termination may be another transmission line having the same characteristic impedance. 
When $Z_L=Z_0$, $\Gamma=0$ and there is no reflection.   

{\bf Open Circuit.} 
\index{open circuit} 
An ``open circuit'' is the absence of a termination.
This condition implies $Z_L\rightarrow\infty$, and subsequently $\Gamma\rightarrow +1$.  
Since the \emph{current} reflection coefficient is $-\Gamma$, the reflected current wave is $180^{\circ}$ out of phase with the incident current wave, making the total current at the open circuit equal to zero, as expected.

{\bf Short Circuit.}  
\index{short circuit} 
``Short circuit'' means $Z_L=0$, and subsequently $\Gamma=-1$.
In this case, the phase of $\Gamma$ is $180^{\circ}$, and therefore, the potential of the reflected wave cancels the potential of the incident wave at the open circuit, making the total potential equal to zero, as it must be.  
Since the \emph{current} reflection coefficient is $-\Gamma=+1$ in this case, the reflected current wave is in phase with the incident current wave, and the magnitude of the total current at the short circuit non-zero as expected.

{\bf Purely Reactive Load.}
\index{load!reactive}
A purely reactive load, including that presented by a capacitor or inductor, has $Z_L = j X$ where $X$ is reactance.  
In particular, an inductor is represented by $X>0$ and a capacitor is represented by $X<0$.
We find
%
\begin{equation}
\Gamma = \frac{-Z_0+j X}{+Z_0+j X}
\end{equation}
%
The numerator and denominator have the same magnitude, so $\left|\Gamma\right|=1$.
Let $\phi$ be the phase of the denominator ($+Z_0+j X$).  
Then, the phase of the numerator is $\pi-\phi$.
Subsequently, the phase of $\Gamma$ is $\left(\pi-\phi\right)-\phi=\pi-2\phi$.
Thus, we see that the phase of $\Gamma$ is no longer limited to be $0^{\circ}$ or $180^{\circ}$, but can be any value in between.
The phase of reflected wave is subsequently shifted by this amount.

{\bf Other Terminations.}
Any other termination, including series and parallel combinations of any number of devices, can be expressed as a value of $Z_L$ which is, in general, complex-valued. 
The associated value of $\left|\Gamma\right|$ is limited to the range 0 to 1.
To see this, note:
%
\begin{equation}
\Gamma = \frac{Z_L-Z_0}{Z_L+Z_0} = \frac{Z_L/Z_0-1}{Z_L/Z_0+1}
\end{equation}
%
Note that the smallest possible value of $\left|\Gamma\right|$ occurs when the numerator is zero; i.e., when $Z_L=Z_0$.
Therefore, the smallest value of $\left|\Gamma\right|$ is zero.
The largest possible value of $\left|\Gamma\right|$ occurs when $Z_L/Z_0\rightarrow\infty$ (i.e., an open circuit) or when $Z_L/Z_0=0$ (a short circuit); the result in either case is $\left|\Gamma\right|=1$.
Thus,
%
%ORB HIGHLIGHT
\begin{mdframed}[backgroundcolor=yellow!20,nobreak=true] 
\begin{equation}
0\le\left|\Gamma\right|\le 1
\end{equation}
\end{mdframed}









